\section*{Chapter \ref{ch:g9:natural-risks} --- Living with Natural Risks}

\begin{solution}{exo:g9:natural-risks:1}
(a) Hazard. (b) \omterm{def:g9:natural-risks:vulnerability}{Vulnerability}. (c) \omterm{def:g9:natural-risks:vulnerability}{Vulnerability}. (d) Hazard. (e) Both:
the mudflows are the hazard; building the town in their path makes it
vulnerable.
\end{solution}

\begin{solution}{exo:g9:natural-risks:2}
(a) False: where no one lives, nothing can be harmed, so the \omterm{def:g9:natural-risks:risk}{risk} is low.
(b) True. (c) True. (d) False: a warning does not change the \omterm{def:g4:volcanoes-earthquakes-first-look:volcano}{volcano}; it
lowers the \omterm{def:g9:natural-risks:vulnerability}{vulnerability} of the people, who can leave in time.
\end{solution}

\begin{solution}{exo:g9:natural-risks:3}
(a) Yes, days ahead. (b) Yes, hours to days ahead. (c) No: \omterm{def:g4:volcanoes-earthquakes-first-look:earthquake}{earthquakes}
cannot be predicted. (d) Yes: once the \omterm{def:g7:earthquakes:tsunami}{tsunami} has started, it takes
hours to cross the \omterm{def:g2:water-on-earth:ocean}{ocean}, and the warning travels faster.
\end{solution}

\begin{solution}{exo:g9:natural-risks:4}
(a) The shaking of the ground: swarms of small \omterm{def:g4:volcanoes-earthquakes-first-look:earthquake}{earthquakes} as \omterm{def:g7:volcanoes:magma}{magma}
rises. (b) The tilt of the slope: the \omterm{def:g4:volcanoes-earthquakes-first-look:volcano}{volcano} swelling. (c) How the
ground moves, by millimetres: the swelling as \omterm{def:g7:volcanoes:magma}{magma} gathers. (d) The
gases of the plume: more gas escaping as \omterm{def:g7:volcanoes:magma}{magma} comes up.
\end{solution}

\begin{solution}{exo:g9:natural-risks:5}
(a) $45 \div 2 = \qty{22.5}{km/h}$. (b) $4 \times 2 = \qty{8}{km}$, far
enough to reach higher ground. (c) Nobody saw the \omterm{def:g4:volcanoes-earthquakes-first-look:eruption}{eruption}, hidden by the
\omterm{def:g5:natural-hazards:storm}{storm}, and nobody warned them: there was no quick way to send the
information to the town, and no evacuation plan. (d) For example: an
evacuation plan with routes to higher ground, practised in advance; a
chain of warning from the scientists to the town; more instruments on the
\omterm{def:g4:volcanoes-earthquakes-first-look:volcano}{volcano}.
\end{solution}

\begin{solution}{exo:g9:natural-risks:6}
(a) The \omterm{def:g5:water-shapes-the-land:flood}{flood} \qty{10}{m} above the \omterm{def:g1:landscapes-around-us:water}{river}. (b) About \qty{0.47}{km} for
the \qty{2}{m} \omterm{def:g5:water-shapes-the-land:flood}{flood}, \qty{0.56}{km} for the \qty{5}{m} \omterm{def:g5:water-shapes-the-land:flood}{flood} and
\qty{0.67}{km} for the \qty{10}{m} \omterm{def:g5:water-shapes-the-land:flood}{flood} (readings within about
\qty{0.05}{km} are fine). (c) Only the \qty{10}{m} \omterm{def:g5:water-shapes-the-land:flood}{flood}: the ground is
higher than 2 and \qty{5}{m} above the \omterm{def:g1:landscapes-around-us:water}{river}. (d) A small rise of the
\omterm{def:g1:landscapes-around-us:water}{river} is needed to reach it, and small \omterm{def:g5:water-shapes-the-land:flood}{floods} come more often than large
ones.
\end{solution}

\begin{solution}{exo:g9:natural-risks:7}
(a) High hazard, low \omterm{def:g9:natural-risks:vulnerability}{vulnerability}: moderate \omterm{def:g9:natural-risks:risk}{risk}. (b) Low hazard, high
\omterm{def:g9:natural-risks:vulnerability}{vulnerability}: moderate \omterm{def:g9:natural-risks:risk}{risk}. (c) High hazard, high \omterm{def:g9:natural-risks:vulnerability}{vulnerability}: high
\omterm{def:g9:natural-risks:risk}{risk}. (d) Low hazard, low \omterm{def:g9:natural-risks:vulnerability}{vulnerability}: low \omterm{def:g9:natural-risks:risk}{risk}.
\end{solution}

\begin{solution}{exo:g9:natural-risks:8}
(a) C, then B, then A. (b) The braces form triangles, which cannot lean
without stretching or squashing a brace: the frame keeps its shape
instead of leaning over. (c) The isolators give way as the ground moves
back and forth beneath them, so much less of the motion reaches the
school.
\end{solution}

\begin{solution}{exo:g9:natural-risks:9}
(a) $1 \div 100 = 1~\%$ per year. (b) The chance is the same every year,
whatever happened the year before: the \omterm{def:g5:water-shapes-the-land:flood}{flood} can come again next year.
(c) $300 \div 100 = 3$ \omterm{def:g5:water-shapes-the-land:flood}{floods}, on average.
\end{solution}

\begin{solution}{exo:g9:natural-risks:10}
(a) $3600 \div 15 = 240$ readings. (b) $240 \times 24 \times 365 =
\num{2102400} \approx 2.1 \times 10^6$ readings. (c) $2400 \div 800 =
\qty{3}{h}$. (d) In these three hours, the warning centres can warn the
\omterm{def:g1:landscapes-around-us:water}{coast} and people can move to higher ground.
\end{solution}

\begin{solution}{exo:g9:natural-risks:11}
(a) $4 \times 3 \times 3.8 = \qty{45.6}{L}$. (b) $45.6 \div 1.5 = 30.4$,
so 31 bottles. (c) For example: food that keeps, to eat when the shops
are shut; a torch, if the electricity is cut; a radio working on
batteries or by hand, to hear the warnings; a first-aid kit, to treat
small injuries; warm blankets. (d) So that it can be taken at once if the
family has to leave quickly.
\end{solution}

\begin{solution}{exo:g9:natural-risks:12}
For example: (1) draw a \omterm{def:g9:natural-risks:hazard-map}{hazard map} of the mudflows, \omterm{def:g1:landscapes-around-us:landform}{valley} by \omterm{def:g1:landscapes-around-us:landform}{valley}
(hazard side); (2) install \omterm{def:g7:earthquakes:seismometer}{seismometers}, tiltmeters, GPS receivers and gas
instruments on the \omterm{def:g4:volcanoes-earthquakes-first-look:volcano}{volcano} (hazard side); (3) set up a fast chain of
warning from the observatory to the town, with sirens (\omterm{def:g9:natural-risks:vulnerability}{vulnerability}
side); (4) plan and signpost routes to higher ground, and practise them
with drills (\omterm{def:g9:natural-risks:vulnerability}{vulnerability} side); (5) build no new homes or schools in the
path of the mudflows (\omterm{def:g9:natural-risks:vulnerability}{vulnerability} side).
\end{solution}

\begin{solution}{pb:g9:natural-risks:1}
\textbf{1.} The \omterm{def:g1:landscapes-around-us:water}{river} is at \qty{143}{m} beside C, \qty{137}{m} beside A
and \qty{126}{m} beside B; on the \omterm{def:g2:drawing-to-map:map}{map} C is to the west and B to the east,
so the \omterm{def:g1:landscapes-around-us:water}{river} flows east.

\textbf{2.} A: $137 - 137 = \qty{0}{m}$; B: $276 - 126 = \qty{150}{m}$;
C: $236 - 143 = \qty{93}{m}$.

\textbf{3.} Only A, which is less than \qty{5}{m} above the \omterm{def:g1:landscapes-around-us:water}{river}.

\textbf{4.} About \qty{0.57}{km} (between 0.5 and \qty{0.65}{km}).

\textbf{5.} The rise over \qty{100}{m} is half the rise over
\qty{200}{m}: A $1.4 \div 2 = 0.7~\%$; B $83 \div 2 = 41.5~\%$;
C $20 \div 2 = 10~\%$.

\textbf{6.} B: its slope is very steep, and loose \omterm{def:g3:stones-pebbles-sand:rock}{rock} and soil can slide
down a steep slope, especially after heavy \omterm{def:g4:clouds-rain-snow:precipitation}{rain}.

\textbf{7.} A: about \qty{1}{km}; B: almost on the \omterm{def:g7:earthquakes:fault}{fault} (less than
\qty{0.05}{km}); C: about \qty{1.3}{km}.

\textbf{8.} If the \omterm{def:g7:earthquakes:fault}{fault} moves, the ground breaks and shifts along it,
which can tear a building apart; and an \omterm{def:g4:volcanoes-earthquakes-first-look:earthquake}{earthquake} on the \omterm{def:g7:earthquakes:fault}{fault} would
shake the ground there very hard.

\textbf{9.}

\begin{center}
{\small
\begin{tabular}{|c|c|c|c|}
\hline
site & in the \omterm{def:g5:water-shapes-the-land:flood}{flood} zone? & steep slope? & on the \omterm{def:g7:earthquakes:fault}{fault}? \\
\hline
A & yes & no & no \\
B & no & yes & yes \\
C & no & no & no \\
\hline
\end{tabular}}
\end{center}

\textbf{10.} Site C: it is the only site with three ``no'': \qty{93}{m}
above the \omterm{def:g1:landscapes-around-us:water}{river}, a gentle slope of \qty{10}{\%}, more than \qty{1}{km}
from the \omterm{def:g7:earthquakes:fault}{fault}.

\textbf{11.} For example: diagonal braces and shear walls in the frame,
strong joints, or isolators under the building; the pupils practise
\omterm{def:g4:volcanoes-earthquakes-first-look:earthquake}{earthquake} drills (drop, cover, hold on).

\textbf{12.} \textbf{17 houses} out of 28 stand inside the \omterm{def:g5:water-shapes-the-land:flood}{flood} zone:
$17 \div 28 \approx 0.61$, about 3 houses in 5.
\end{solution}
